\documentclass[12pt]{article}
\usepackage{amsmath,amssymb}

\begin{document}
\section*{{2019 AIME II Problems/Problem 1}}
Source: AoPS Wiki

\subsection*{Solution}
[asy] unitsize(10); pair A = (0,0); pair B = (9,0); pair C = (15,8); pair D = (-6,8); pair E = (-6,0); draw(A--B--C--cycle); draw(B--D--A); label("$A$",A,dir(-120)); label("$B$",B,dir(-60)); label("$C$",C,dir(60)); label("$D$",D,dir(120)); label("$E$",E,dir(-135)); label("$9$",(A+B)/2,dir(-90)); label("$10$",(D+A)/2,dir(-150)); label("$10$",(C+B)/2,dir(-30)); label("$17$",(D+B)/2,dir(60)); label("$17$",(A+C)/2,dir(120)); draw(D--E--A,dotted); label("$8$",(D+E)/2,dir(180)); label("$6$",(A+E)/2,dir(-90)); [/asy] - Diagram by Brendanb4321 Extend $AB$ to form a right triangle with legs $6$ and $8$ such that $AD$ is the hypotenuse and connect the points $CD$ so
that you have a rectangle. (We know that $\triangle ADE$ is a $6-8-10$ , since $\triangle DEB$ is an $8-15-17$ .) The base $CD$ of the rectangle will be $9+6+6=21$ . Now, let $O$ be the intersection of $BD$ and $AC$ . This means that $\triangle ABO$ and $\triangle DCO$ are similar with ratio $\frac{21}{9}=\frac73$ . Set up a proportion, knowing that the two heights add up to 8. We will let $y$ be the height from $O$ to $DC$ , and $x$ be the height of $\triangle ABO$ . \[\frac{7}{3}=\frac{y}{x}\] \[\frac{7}{3}=\frac{8-x}{x}\] \[7x=24-3x\] \[10x=24\] \[x=\frac{12}{5}\] This means that the area is $A=\tfrac{1}{2}(9)(\tfrac{12}{5})=\tfrac{54}{5}$ . This gets us $54+5=\boxed{059}.$ -Solution by the Math Wizard, Number Magician of the Second Order, Head of the Council of the Geometers

\end{document}
